\begin{lexlean}{IdsAU}
\useglossary{lexlean.std.nat@1.2.0}
\useglossary{spec.titles@1.0.0}
\importmodule{Rows}
\title{Authoring conformance claims}
\begin{semanticmodule}
definition{
  name: au01,
  parameters: [],
  result: Rows::IdRow,
  body: (record Rows::IdRow {id := "AU-01", level := "build", suite := "authoring", statement := "Every value of the closed semantic-module schema has one notation and the notation reads back to the value: every alternative of the declaration, term, type, and proof unions, with and without its optional members, is written and read back equal, a node that no readable form carries exactly is written generically, the notation is its own canonical form, and every committed semantic module is written and read back equal to its data.", sections := "§17.21"}),
  executable: true
}

definition{
  name: au02,
  parameters: [],
  result: Rows::IdRow,
  body: (record Rows::IdRow {id := "AU-02", level := "build", suite := "authoring", statement := "A language-1.2 semantic module written in the authoring notation builds, with source maps that relate each declaration to its text in the notation; fmt leaves the notation unchanged and restores its canonical layout from another layout, and the notation under language 1.1 is refused with LLP2003; its data form is not a language-1.2 source (AU-13).", sections := "§17.21"}),
  executable: true
}

definition{
  name: au03,
  parameters: [],
  result: Rows::IdRow,
  body: (record Rows::IdRow {id := "AU-03", level := "build", suite := "authoring", statement := "Malformed notation fails closed and never panics: an unclosed delimiter, a keyword used as a name, a literal that is not canonical, an unknown kind, and every single-token deletion, duplication, and replacement of a committed module fail with a registered code at a position in the notation or link to a module that decodes, and a numeral with a redundant leading zero is LLL1003 outside a string and is not one inside it.", sections := "§17.21"}),
  executable: true
}

definition{
  name: au04,
  parameters: [],
  result: Rows::IdRow,
  body: (record Rows::IdRow {id := "AU-04", level := "build", suite := "authoring", statement := "No language-1.2 source of the product is written as data: every committed language-1.2 source with a semantic, tooling, or core module, in the compiler, the language libraries, the examples, and the positive fixtures, is written in the authoring notation and is canonical, the data form remains only under language 1.0 and 1.1 and in refused fixtures whose payload the notation cannot spell, and a refused fixture whose payload the notation can spell is a defect.", sections := "§17.21"}),
  executable: true
}

definition{
  name: au05,
  parameters: [],
  result: Rows::IdRow,
  body: (record Rows::IdRow {id := "AU-05", level := "build", suite := "authoring", statement := "A language-1.2 tooling module written in the authoring notation renders, with source maps that relate each declaration to its text in the notation; the notation and the canonical JSON of every committed tooling module are in bijection and the notation is its own canonical form; fmt leaves the notation unchanged; and malformed notation, and every single-token deletion and duplication of a committed tooling module, fails with a position or decodes, and nothing panics.", sections := "§17.20, §17.21"}),
  executable: true
}

definition{
  name: au06,
  parameters: [],
  result: Rows::IdRow,
  body: (record Rows::IdRow {id := "AU-06", level := "build", suite := "authoring", statement := "The compiler-semantics ID does not depend on the form a module of the embedded tree is written in: a semantic, tooling, or core module in the authoring notation contributes its data form, so the committed tree and the same tree with every such module written as data have one ID, a change to a module's meaning changes the ID in either form, the identity a preservation record states for the native library is over the same data forms, so it is the committed one, and the embedded tree holds at least one module written in the notation.", sections := "§17.21, §21.2"}),
  executable: true
}

definition{
  name: au07,
  parameters: [],
  result: Rows::IdRow,
  body: (record Rows::IdRow {id := "AU-07", level := "build", suite := "authoring", statement := "No product source is the output of a generator: every `.lex.tex` source of `compiler/`, `language/`, and `examples/` is written and edited as source, and no writer of the repository (the generators `cargo xtask check-calculus` compares and rewrites) emits one of them except the four fixture modules of `compiler/` (`TargetFixtures`, `ReasoningFixtures`, `GnafFixtures`, `GnafCertificates`), generated from committed fixtures, requests, and certificates, and the extraction adapter's module, generated from the registry the pinned Lean answers; a source emitted by any other writer is a defect.", sections := "§17.14, §17.21"}),
  executable: true
}

definition{
  name: au08,
  parameters: [],
  result: Rows::IdRow,
  body: (record Rows::IdRow {id := "AU-08", level := "build", suite := "authoring", statement := "The authoring API reads, writes, converts, and assembles the source of a module without a project: for every committed semantic and tooling source and every core source of at most twelve megabytes it reads the module it links to, writes it back in the form it has and in the other form, assembles the source from its header and data to the canonical notation, and is idempotent in each; and a source that is not a module, or whose notation is malformed, is an error with a position in the source and never a panic.", sections := "§24.8"}),
  executable: true
}

definition{
  name: au09,
  parameters: [],
  result: Rows::IdRow,
  body: (record Rows::IdRow {id := "AU-09", level := "build", suite := "authoring", statement := "Every core module of the closed schema has one notation and the notation reads back to the module and to its data form: every committed core module reads back equal from its notation, its canonical JSON is the canonical JSON the compiler accepts and is, for a module written as data, byte for byte the data it was read from, the notation is its own canonical form, a node is written inside the node that uses it exactly when it is used once by the node after it and every other node is a statement, and a module that carries every kind of node, level, annotation, and declaration member is written and read back equal.", sections := "§17.21.8"}),
  executable: true
}

definition{
  name: au10,
  parameters: [],
  result: Rows::IdRow,
  body: (record Rows::IdRow {id := "AU-10", level := "build", suite := "authoring", statement := "A language-1.2 core module written in the authoring notation builds, with source maps that name the notation; fmt leaves the notation unchanged and restores its canonical layout from another layout; and the notation under language 1.0 and 1.1 is refused with LLP2003; its data form is not a language-1.2 source (AU-13).", sections := "§17.21.8"}),
  executable: true
}

definition{
  name: au11,
  parameters: [],
  result: Rows::IdRow,
  body: (record Rows::IdRow {id := "AU-11", level := "build", suite := "authoring", statement := "Malformed core notation fails closed and never panics: an unclosed delimiter, a keyword used as a name, a label that is not the number of its node, a forward reference to a node, a numeral with a redundant leading zero, a statement that defines no node, a nesting beyond the bound, and every single-token deletion and duplication of a committed module fail with a registered code at a position in the notation or link to a module that decodes.", sections := "§17.21.8"}),
  executable: true
}

definition{
  name: au12,
  parameters: [],
  result: Rows::IdRow,
  body: (record Rows::IdRow {id := "AU-12", level := "build", suite := "authoring", statement := "A core module written with names for its statements and its variables reads to exactly the module its canonical notation reads to, and the canonical notation keeps the names of the binders and writes the numbers of the statements: a variable is written by the name of its binder exactly when no nearer binder has the name, a constant a binder in scope is named like is written as a string, a theorem added to the certification library by name reads to the committed module and is written as the committed bytes, and an unknown or repeated statement name and a variable given universe levels are refused with a position.", sections := "§17.21.8"}),
  executable: true
}

definition{
  name: au13,
  parameters: [],
  result: Rows::IdRow,
  body: (record Rows::IdRow {id := "AU-13", level := "build", suite := "authoring", statement := "Under language 1.2 a module written with a data argument (`\\semanticdata{…}`, `\\toolingdata{…}`, or `\\coredata{…}`) is refused with LLP2003 at the position of the data argument and naming the authoring notation, whichever of the three environments holds it; fmt does not rewrite it and reports the same refusal; lexlean::sdk::to_data_form is the one conversion into the data form and to_notation the one out of it; and the data form is still accepted under language 1.1, where the semantic example links, and is still the form of the Atlas under language 1.0.", sections := "§17.11, §17.21, §24.8"}),
  executable: true
}
\end{semanticmodule}
\end{lexlean}
